\section{Introduction}
\label{sec:intro}

A chemist optimizing a cross-coupling chooses a ligand, a base and a solvent, and sometimes a temperature and a concentration. Most of this search space is categorical, and each experiment costs hours of instrument time. Bayesian optimization (BO) handles small budgets well when its inputs carry a useful notion of distance \citep{shields2021bayesian}. For categorical inputs that distance has to come from somewhere. One-hot encodings provide none. DFT or molecular descriptors provide one, but a chemist must choose and compute them again for every reaction class.

Large language models (LLMs) offer a second source of distance, since they have read the literature on phosphine ligands and carbonate bases. Recent work draws on that knowledge in three ways. One line embeds text descriptions of experiments and fits a Gaussian process (GP) on the embeddings, with or without finetuning the LLM \citep{kristiadi2024sober,rankovic2026gollum}. A second lets the LLM act as the surrogate model or as the optimizer itself \citep{liu2024llambo,ramos2026boicl,macknight2025pretrained}. A third turns LLM output into a prior over the objective \citep{yuan2026lgbo,chen2026caliper}. All three ask the model for an absolute quantity: a vector, a predicted yield, or a score for each candidate.

Absolute quantities are where LLM output is least dependable. LLM surrogates shift their predictions with the wording and order of the prompt \citep{lei2026elicitation}, self-reported confidence can rise with error \citep{chen2026caliper}, and LLM optimizers re-propose experiments they have already run at a far higher cost per iteration than a GP \citep{macknight2025pretrained}. Much of what a chemist knows about reagents, by contrast, is relational: Cs$_2$CO$_3$ is a stronger base than K$_2$CO$_3$, and P($t$Bu)$_3$ is bulkier than PPh$_3$, whatever their exact descriptor values. Psychometrics has preferred paired comparisons to absolute ratings since \citet{thurstone1927law}, and work on LLM judges reports the same preference for text evaluation and ranking \citep{liusie2024comparative,qin2024prp}.

We test whether this carries over to reagent properties. On four reaction datasets where DFT descriptors provide a ground truth, we ask five LLMs which of two reagents has more of a named property and, separately, to rate each reagent from 1 to 10. Comparisons agree with the descriptor ordering for \PairAccPct{} of pairs on average, orderings implied by the same models' ratings for \RateAccPct{}, and in-context yield predictions correlate with measured yields at Spearman $\rho=\PredRho$ (Figure~\ref{fig:teaser}). The gap appears for every model we tried.

Turning ordinal answers into geometry for BO requires a model, and that model has to stop trusting the LLM when the experiments disagree with it. We introduce \ours{}, a BO method built around that one question. It places every reagent at a point in a low-dimensional latent space (Figure~\ref{fig:method}). The LLM names a few axes, such as \emph{ligand steric bulk}; the remaining axes stay unnamed and are fit to the data. A GP models the objective on these coordinates. An ordered-probit likelihood links each LLM answer to the coordinates through a reliability $\lambda_k$ that is inferred jointly with them, so the experiments decide how much each named property counts. Between experiments, \ours{} asks only the comparisons whose answers could change the next experiment. With every $\lambda_k$ at zero, \ours{} is a latent-variable GP \citep{zhang2020latent}, so a misleading LLM costs at most the information its answers would have added.

Table~\ref{tab:robust} shows what this buys. Among the LLM-based methods we evaluate, \ours{} has the lowest regret when the LLM's chemistry is right, and it is the only one that stays near a knowledge-free GP when we permute the outcomes so that the LLM's chemistry is wrong. Our contributions are (i) a measurement of comparative versus absolute LLM judgments of reagent properties across five models (\S\ref{sec:analysis}); (ii) \ours{}, which infers reagent coordinates jointly from experiments and comparisons, learns a reliability per named axis, chooses comparisons by their effect on the next decision, and names new axes when unnamed dimensions start to matter (\S\ref{sec:method}); (iii) an analysis showing that \ours{} contains the latent-variable GP, that its reliabilities are identifiable once named axes are standardized, and that its input dimension stays between 8 and 14 (\S\ref{sec:props}); and (iv) an evaluation on nine reaction optimization tasks against fourteen baselines, with oracles that are made wrong, inverted or stripped of context (\S\ref{sec:exp}).
